% !TeX TS-program = xelatex
% 虚构演示简历 — 仅用于 hellojobs 测试数据
\documentclass{resume-photo}
\usepackage{xeCJK}
\setCJKmainfont{Noto Serif CJK SC}
\usepackage{fontspec}
\usepackage{enumitem}
\setlist[itemize]{itemsep=0pt, topsep=1pt, parsep=0pt, partopsep=0pt}

\ResumeName{彭宇}

\begin{document}

\ResumeContacts{
  138-0027-0027,%
  \ResumeUrl{mailto:pengyu@ucas.ac.cn}{pengyu@ucas.ac.cn},%
  \textnormal{中国科学院大学 | 网络空间安全 · 硕士 | 2023—2026}%
}

\ResumeTitle

\noindent 隐私计算与联邦学习，熟悉 PySyft 思想、秘密分享与安全聚合；横向联邦图像分类在 Non-IID 下准确率损失 <2\%。注重可观测性与文档沉淀，习惯在评审阶段拆解风险；参与线上复盘与跨团队联调，英语 CET-6，可阅读官方文档与 RFC（演示数据）。

\section{教育背景}
\ResumeItem
{中国科学院大学}
[\textnormal{网络空间安全学院 | 网络空间安全 | 硕士}]
[2023.09—2026.06(预计)]

\begin{itemize}
  \item 研究所：信息工程研究所（虚构组名略）。
\end{itemize}


\ResumeItem
{企业开放日与实训}
[\textnormal{短期实训营（演示）}]
[2023—2024]

\begin{itemize}
  \item 参加头部互联网公司 2 周封闭实训，完成从需求到上线的完整小组项目。
  \item 在实训中负责接口设计与联调，小组答辩成绩排名前 20\%。
  \item 获得实训营「优秀学员」与内推绿色通道（测试数据，勿当真）。
  \item 沉淀实训笔记 1.5 万字，含架构图与排错记录。
\end{itemize}



\section{专业技能}
\begin{itemize}
  \item 熟悉 Python、PyTorch、差分隐私基础与 DP-SGD。
  \item 掌握安全聚合（SecAgg）与同态加密入门实验。
  \item 了解 MPC 协议仿真与通信轮次分析。
  \item 熟悉 Linux 日常运维与 Shell，具备日志检索与 CPU/内存突增初步定位能力。
  \item 掌握 Git / CI 与 Code Review，了解 Docker 与 Kubernetes 基本编排与滚动发布。
  \item 了解 OAuth2 / JWT、接口幂等与 REST 规范，具备单元测试与集成测试习惯（JUnit/pytest 等）。
  \item 能阅读英文技术文档，具备跨团队沟通与需求澄清能力（测试简历）。
\end{itemize}

\section{实习经历}
\ResumeItem
{蚂蚁集团 | 隐私计算}
[\textnormal{研究型实习生}]
[2024.06—2024.12]

\begin{itemize}
  \item \textbf{职责}: 纵向联邦特征对齐模块原型验证。
  \item \textbf{成果}: PSI 匹配耗时在 100 万级 ID 场景 <90s（内网）。
\end{itemize}


\ResumeItem
{海康威视 | 行业应用}
[\textnormal{嵌入式 / 后端实习生}]
[2023.09—2024.02]

\begin{itemize}
  \item \textbf{设备接入}: 参与国标设备接入网关协议适配与异常码映射。
  \item \textbf{可靠性}: 实现断线重连与心跳退避策略，掉线恢复时间缩短 40\%。
  \item \textbf{日志}: 统一日志格式与 traceId 透传，问题定位时间减半。
  \item \textbf{客户}: 支持 2 次驻场联调，收集需求并转化为可验收用例。
  \item \textbf{规范}: 熟悉编码规范与静态检查门禁，CI 全绿方可合并。
\end{itemize}



\section{项目经历}
\ResumeItem
{联邦学习仿真平台}
[\textnormal{实验室项目}]
[2024.02—2024.08]

\begin{itemize}
  \item 支持 10 客户端异步聚合，收敛曲线与中心化基线对齐度 0.93。
\end{itemize}


\ResumeItem
{特征存储与在线推理}
[\textnormal{Feast 类实践（演示）}]
[2023.08—2024.01]

\begin{itemize}
  \item \textbf{一致}: 离线在线特征对齐，监控 PSI 漂移。
  \item \textbf{延迟}: 在线特征查询 P99 <6ms，支撑排序模型实时更新。
  \item \textbf{平台}: 与模型服务 gRPC 联调，完成金丝雀发布。
  \item \textbf{文档}: 输出特征字典与血缘说明，供算法与工程共用。
\end{itemize}



\section{个人荣誉}
\begin{itemize}
  \item 中国科学院大学 三好学生
  \item 密码创新竞赛 优胜奖
  \item 全国计算机等级考试四级（网络工程师）（演示）
  \item 校级「优秀学生干部 / 优秀团员」或技术社团骨干（综合测评前 15\%）
  \item 开源贡献：向知名项目提交被合并的 PR（文档/测试/feature）（演示）
\end{itemize}

\end{document}
